\chapter{Reading nested data}
\MARSChapterHeading{Reading nested data}
\label{revision:cbor}

A sensor message may contain a counter, a binary device identifier and several
signed readings. A fixed record could put each field at a predetermined offset.
A structured message instead carries descriptions of its values. The readings
can be an array, and an optional field can contain another map. Finding the next
field requires understanding where that whole nested value ends.

\href{https://www.rfc-editor.org/rfc/rfc8949.html}{CBOR}, the Concise Binary Object
Representation, supplies those descriptions. Our worked program accepts a bounded
subset and projects a telemetry message into a flat application record. It then
writes a deterministic encoding of the required fields. The program is complete
in \texttt{revision/labs/files/cbor-telemetry.asm}; its caller selects the head
reader, walker, decoder, size calculator, encoder or their combined pipeline.

This chapter assumes unsigned bounds, byte order, bit fields, nested function
calls and validation before publication. Its new problem is retaining a parent's
unfinished work while visiting children. The independent exercise changes the
outer framing: it recognizes an entire sequence of supported items instead of
one telemetry record.

\section{Separate a value's description from its contents}

Each CBOR item starts with a \emph{head}. The initial byte divides into a major
type in its upper three bits and additional information in its lower five.
For this exercise, types zero and one describe nonnegative and negative integers,
two a byte string, four an array, and five a map. A byte string is counted binary
data; neither a zero byte nor a byte resembling a head ends its payload.

For \texttt{82}, shifting right five produces type four, and masking with 31
produces argument two. This is an array containing two items. For \texttt{42},
the same argument instead counts two payload bytes because the type is two.
An argument acquires its unit from the type. A map argument counts pairs, so a
map of two pairs still requires two keys and two values.

\begin{table}[htbp]
\centering
\begin{tabular}{r r r l}
\toprule
Low five bits & Following bytes & Head bytes & Interpretation\\
\midrule
0--23 & 0 & 1 & Argument is the low field\\
24 & 1 & 2 & Next byte is the argument\\
25 & 2 & 3 & Next two bytes, big-endian\\
26 & 4 & 5 & Next four bytes, big-endian\\
27 & 8 & 9 & Next eight bytes, big-endian\\
\bottomrule
\end{tabular}
\caption{Supported head widths. A width selector is not itself the argument. Low fields 28--30 are reserved; indefinite-length form 31 is outside this exercise's definite-length subset.}
\label{tab:revision-cbor-widths}
\end{table}


Low field 24 requests a following argument byte. Thus \texttt{18 17} represents
unsigned 23, while \texttt{18 18} represents unsigned 24. Field 25 requests two
bytes; \texttt{19 01 00} represents 256. The head reader accepts supported arguments
in wider-than-necessary heads. Later, the writer chooses their shortest form.
Acceptance and preferred output spelling are separate decisions.

The full-width counter head is \texttt{1b 01 02 03 04 05 06 07 08}. Establish
that all nine head bytes are available before reading its eight argument bytes.
Starting with zero, shift the accumulator left eight and combine each unsigned
byte. The result is \texttt{0102030405060708}. Use a doubleword shift:
\texttt{dsll} retains the full 64-bit accumulation; a word shift would discard
part of this counter.

\MARSCodeLabel{Accumulate the argument without narrowing it}
\begin{minted}{mars-idea}
    daddiu $t5, $t5, 1
    lbu $t0, 0($t5)
    dsll $t6, $t6, 8
    or $t6, $t6, $t0
    addiu $t3, $t3, -1
    bnezc $t3, head_bytes
    nop
\end{minted}

All counter bit patterns are valid, including eight
\texttt{ff} bytes, regardless of a debugger's signed interpretation.

\texttt{cbor\_head} receives byte pointer and unsigned available count in
\texttt{a0/a1}. Successful \texttt{v1} counts head bytes; \texttt{a2} is the type
and \texttt{a3} the unsigned argument bits. A successful head establishes neither
a byte-string payload nor any array children. On failure, status is one and all
three result registers are zero. Zero available bytes fail before pointer access;
a large count does not make an unsupported type valid.

\section{Keep one continuation for every unfinished container}

Consider the optional value in the worked message:
\texttt{a2 01 82 42 80 ff 21 00 80}. It is a two-pair map. Key one maps to an
array of two values: byte string \texttt{80 ff}, then integer minus two, encoded
\texttt{21}. Key zero maps to an empty array, \texttt{80}. The byte \texttt{ff}
inside the string is payload. Treating it as a separate head would lose the
boundary established by \texttt{42}.

\begin{figure}[htbp]
\centering
\begin{tabular}{l r r l}
\toprule
Active container & Cursor offset & Outstanding work & Retained identity\\
\midrule
Outer telemetry map & 3 & Value for key 3; three pairs follow & Key 3 seen\\
Optional map & 5 & Value for key 1; one pair follows & Key 1 seen\\
Array & 6 & Two values & No key mask\\
\bottomrule
\end{tabular}
\caption{Suspended work while entering the byte string at offset six in the trace. The array head was at five. The array pauses at child offset six. Reading that byte string's head advances the child's cursor to seven; its payload ends at nine. Each parent's cursor and outstanding work must survive that child. These are logical obligations, not required scratch registers or a prescribed frame layout.}
\label{fig:revision-cbor-parent-state}
\end{figure}


A walker keeps a cursor and remaining byte budget. A primitive integer ends
after its head. A byte string additionally consumes its checked payload count.
An array consumes its declared number of child items. A map consumes a checked
key and one complete child value per pair. After a child succeeds, advance by
its returned byte extent and reduce both the byte budget and the parent's
outstanding child or pair count.

The worked walker uses recursive calls. Before calling the child, it saves the
parent cursor, remaining bytes, depth, outstanding work and map mask in its own
frame. A second implementation can keep these continuations in an explicit
stack and run one loop. Both must retain the same information. Merely saving the
return address preserves where instructions resume; it does not preserve the
parent's data.

One particularly easy mistake is leaving depth in \texttt{a2} across a call to
\texttt{cbor\_head}. That helper returns the major type in \texttt{a2}. Save the
inherited depth before the call; otherwise a type code becomes a nesting count.
The provided recursive walker keeps depth in a saved register belonging to the
current frame. Each recursive child preserves its caller's full saved values.

\section{Apply local bounds to every nested value}

Our application permits definite integers, byte strings of at most 32 bytes,
arrays of at most 64 items, and maps of at most eight pairs. Every map uses
unique unsigned keys zero through seven. These rules apply inside optional
metadata too. Text, tags, floating-point and simple values, and indefinite
lengths are outside this subset. Several are valid CBOR elsewhere; rejecting
them describes this exercise's profile, not a defect in the general format.

Track seen map keys with an eight-bit mask. For key $k$, form $1\ll k$, reject
if that bit is already set, and then set it. Compare decoded keys: \texttt{00}
and \texttt{18 00} name the same key zero. \MARSCodeLabel{Check decoded key identity in the current map}
\begin{minted}{mars-idea}
    li $t0, 1
    sllv $t0, $t0, $a3
    and $t1, $s5, $t0
    bnezc $t1, skip_bad
    nop
    or $s5, $s5, $t0
\end{minted}

A nested map starts its own mask.
Reusing the parent's mask would reject a valid child key; forgetting the mask
would accept a repeated field with ambiguous ownership.

At most eight containers may be open simultaneously. The telemetry root is
one. Primitives do not add depth. With inherited depth seven, an empty array
can open the eighth container and close immediately. At inherited depth eight,
even an empty array would open a ninth and must fail. A sequence of sibling
containers does not accumulate depth: each sibling starts with the same
inherited count after its predecessor closes.

\texttt{cbor\_skip} takes pointer, unsigned byte count and inherited depth in
\texttt{a0/a1/a2}. It first rejects count above 1024 with status two, then depth
above eight with status one, before pointer access. Zero count fails with one.
Success returns the byte extent of \emph{one} supported item. The worked optional
map begins at offset three and occupies nine bytes. The whole root occupies 52.
A following byte does not enlarge either item. The telemetry decoder, rather
than the generic walker, must establish that the selected input ends there.

\section{Turn one complete map into an application record}

The telemetry schema requires keys zero, one and two exactly once. Their values
are respectively an unsigned 64-bit counter, a binary identifier of at most 32
bytes, and an array of at most 64 signed 16-bit readings. Optional keys three
through seven carry supported nested values. They must validate even though the
flat output will not retain them. Dropping a field does not authorize ignoring
its malformed contents.

The 52-byte trace deliberately places optional key three first, readings second,
identifier third and counter last. Its wider root head \texttt{b8 04} declares
four pairs. After the nine-byte optional value, key two at offset twelve precedes
the ten-reading array. Key one at 35 precedes identifier \texttt{00 ff 41}.
The wider key-zero head at 41 precedes the full-width counter at 43. Required
fields need not arrive in output order.

\begin{table}[htbp]
  \centering
  \begin{tabular}{r l l}
    \toprule
    Byte offset & Local field                       & Active extent                  \\
    \midrule
    0           & Unsigned 64-bit counter           & Eight bytes                    \\
    8           & Unsigned 32-bit identifier length & Four bytes                     \\
    12          & Unsigned 32-bit reading count     & Four bytes                     \\
    16          & Binary identifier                 & Length bytes, at most 32       \\
    48          & Signed 16-bit readings            & Twice count bytes, at most 128 \\
    \bottomrule
  \end{tabular}
  \caption{The eight-byte-aligned local record reserves 176 bytes. Its integers use native little-endian storage. Successful decoding writes the prefix and active payloads; unused identifier and sample tails retain their previous bytes. No field points back into the input.}
  \label{tab:revision-cbor-record}
\end{table}


CBOR negative integers use argument $u$ to mean $-1-u$. Thus \texttt{20} means
minus one and \texttt{37} means minus 24. The head \texttt{38 18} carries
argument 24 and means minus 25. To accept signed 16-bit readings, either
integer type needs argument at most 32767. Type zero then gives 0 through 32767;
type one gives minus one through minus 32768. Negating the argument alone would
make every negative reading one too large.

Keep the output private until the root, required-key mask and exact ending have
all passed. The reference builds a candidate in its frame, then publishes only
the prefix and active payloads. An alternative remembers checked input slices
and revisits them after validation. Under the disjoint, read-only input contract,
those bytes remain stable for the second pass. Neither strategy may expose a
plausible identifier before discovering that the final counter is truncated.

\texttt{telemetry\_decode} receives input, unsigned complete count and an
eight-byte-aligned record pointer in \texttt{a0/a1/a2}. Count above 1024 fails
with status two, and count below seven with one, before reading either pointer.
Success returns the complete input count in \texttt{v1}. Failure returns zero
there and leaves the whole caller record unchanged. A successfully walked root
followed by another item is still a failed telemetry input.

\section{Derive an exact size before writing}

The flat record no longer contains optional metadata or original head widths.
Encoding it is an application projection. It cannot reproduce the original map
byte for byte, and it does not promise to do so. The writer emits required keys
in order zero, one, two with shortest heads. For these small unsigned keys,
that order agrees with the core deterministic ordering in RFC 8949.

Let $H(u)$ count the shortest head bytes for unsigned argument $u$: one, two,
three, five or nine according to the width table. Let $A(x)=x$ for a nonnegative
reading and $A(x)=-1-x$ otherwise. For counter $T$, identifier length $L$ and
reading count $N$, the exact encoded size is
\[
4+H(T)+H(L)+L+H(N)+\sum_{i=0}^{N-1}H(A(x_i)).
\]
The four fixed bytes are the root map head and three keys. Identifier contents
do not affect the size; sample values do. The bound function therefore reads
the sixteen-byte prefix and active samples, but no identifier payload.

For the trace the terms are $4+9+1+3+1+19=37$. A zero counter with empty identifier
and empty readings needs seven bytes. The maximum result is 241 bytes. Unlike
QOI's conservative reservation contract, this encoder must accept the exact
37-byte trace buffer; demanding 241 bytes would reject sufficient capacity.

Both bound and encoder reject identifier length above 32 or reading count above
64 before touching payload storage. Encoding then computes the size and checks
capacity unsigned before reading identifier bytes or writing output. Once this
preflight passes, every selected payload has a proved destination extent.
Successful \texttt{v1} is the exact byte count. No terminator belongs after it.

\section{Distinguish syntax, schema and projection when testing}

A head reader, generic walker and telemetry decoder answer different questions.
The head \texttt{98 40} is complete even without its 64 array elements. A complete
array can be supported by the walker while failing the telemetry schema because
the root must be a map. A valid telemetry map can have several accepted byte
representations while projecting to one preferred required-field encoding.
Use a separate counterexample for each boundary.

Run mode five of the supplied caller to join head reading, walking, decoding,
bounding and encoding. On the trace, walking and decoding each account for all
52 input bytes, while the encoded output occupies 37. Compare decoded field
values and required output bytes, not an expectation that the two byte counts
must match. Inspect the binary identifier as three counted bytes; a console
string routine would stop at its first zero.

All five entries preserve the full saved registers, frame and stack pointers,
and return through the incoming return address. They do no console or file I/O.
Every completed failure clears \texttt{v1}; status one means unsupported or invalid
structure, field or depth, and two means excessive selected input or insufficient
output capacity as specified for that entry. Byte pointers may be unaligned;
caller objects are disjoint. The lab's sentinels and access checks distinguish
active fields from untouched tails and successful publication from early stores.

\section{Recognize a complete sequence independently}

Complete \texttt{cbor\_sequence\_count} in \texttt{cbor-sequence.asm}. It receives
byte pointer, unsigned complete count and inherited depth in \texttt{a0/a1/a2}.
Accept a concatenation of zero or more supported items under the same local
value rules. Return status zero and the number of top-level items in \texttt{v1}.
This is an item count, not a byte count or the number of descendants.

Reject count above 1024 with status two before checking depth or reading the
pointer. Then reject inherited depth above eight with status one. With a valid
depth, an empty selected sequence succeeds with count zero without pointer
access. Any malformed item, unsupported value or incomplete final suffix fails
with one and clears \texttt{v1}, even after several valid items. Preserve all
caller bytes and the same saved-register contract. Perform no I/O.

For \texttt{00 82 01 02 40}, the result is three: integer zero, one two-element
array, and one empty byte string. The array's two children do not add top-level
items. Appending \texttt{1b} makes the entire sequence invalid because its final
head requests eight missing bytes. A valid-prefix count is not the result on
failure.

The supplied walker may be reused. Keep the accumulated item count and original
inherited depth across each call, advance by the successful byte extent, and
require the whole selected interval to be consumed. An implementation using an
end address can be equivalent to one using a remaining count. Explain why siblings
reuse depth, why every successful iteration advances, and why a late failure
must not return the accumulated prefix count. These arguments should guide your
tests before you inspect the public feedback.
